\section{Cardinal Against His Will (1599--1602)}

\subsection{``Because the Church of God had not his equal in learning''}

On 3~March 1599 Clement~VIII created Bellarmine cardinal priest of
the title of Santa Maria in Via. Two forms of the pope's stated
reason survive in the tradition. The 1907 Catholic Encyclopedia
gives it as: ``the Church of God had not his equal in learning.''
The canonization decretal quotes it in Latin: \latin{``Hunc eligimus
quia non habet parem Ecclesia Dei quod ad doctrinam''}---``we
choose this man because the Church of God has not his equal in
doctrine''---and adds that the pope acted on a man \latin{invitum
et frustra reluctantem}, ``unwilling and resisting in vain''
(\work{Lux illa}, AAS~22, 595; M, resting on process testimony).
The reluctance is not merely hagiographic varnish; it is of a piece
with everything the autobiography says about honors, and with the
conclave prayer reported below. What the new cardinal did with the
income is also on record at the same evidentiary level: he refused
pensions offered by the king of Spain, accepted from the pope only
the modest sum needed to sustain the dignity, and settled on three
rules---to keep the Society's austerity of table and prayer, to
enrich no relatives, and to give to churches or the poor whatever
income remained (\work{Lux illa}, AAS~22, 595; autobiography,
D\"ollinger--Reusch, 41; A/M).

\subsection{The Holy Office, and the Bruno case stated at its ceiling}

As consultor from c.~1597 and as cardinal inquisitor from 1599,
Bellarmine sat inside the two Roman congregations---Holy Office and
Index---whose acts have most darkened his modern reputation. This
study states his documented share exactly and refuses to state
more. In the case of Giordano Bruno, burned in the Campo de' Fiori
on 17~February 1600: the 1908 Catholic Encyclopedia article on
Bruno records the Roman trial's length, the sentence of January
1600, the handing over to the secular arm on 8~February, and the
execution, and states that Bruno ``was not condemned for his
defence of the Copernican system of astronomy, nor for his doctrine
of the plurality of inhabited worlds, but for his theological
errors,'' of which it gives examples (K/E, reported); it does not
name Bellarmine. That Bellarmine belonged to the Inquisition in
those years is certain from its own documents (he appears among
the cardinal inquisitors general in the printed Index decree of
1616: Opere XIX, 322--323; N). The specific claims familiar from
modern accounts---that he framed the list of eight propositions
put to Bruno, that he pressed for or against particular
outcomes---rest on the surviving summary of the lost trial dossier
and on modern reconstructions of it, neither of which was inspected
for this study. They are therefore retained here as uninspected
leads, neither asserted nor denied.

\witnesslimit{Ceiling statement: Bellarmine was one of the cardinal
inquisitors during Bruno's Roman trial and execution; the
institution he served condemned Bruno to death. His individual
votes, arguments, and causal weight are not documented in any
witness read for this study. Neither his sanctity nor his office
settles the question, and no allocation of personal responsibility
is made here in either direction. What his own theology taught
about the punishment of heretics is a separate, documentable
matter, discussed with the reception history
(\S\ref{sec:reception}).}

\subsection{\latin{De auxiliis}: the dispute that was not decided}

The bitterest purely theological quarrel of the age---between the
Dominican theology of physical premotion and the Jesuit theology of
\latin{scientia media}, both claiming to reconcile efficacious grace
with free will---had been referred to Rome in 1597--98. The 1908
Catholic Encyclopedia article on the \latin{Congregatio de auxiliis}
gives the frame: a commission under Cardinals Madruzzi and Aragone
from January 1598; disputations before Clement~VIII in person from
1602 to 1605 (sixty-eight sessions); seventeen more under Paul~V
(K/E, reported). Bellarmine's role was double. As the most
authoritative Jesuit theologian in Rome he was inevitably a
partisan---his own fourth tome \work{De gratia} stood near the
center of the dispute---and after 1599, as cardinal, he attended
sessions. But his documented advice ran against deciding at all:
the 1907 Catholic Encyclopedia reports his counsel that ``the
doctrinal question should not be decided authoritatively, but left
over for further discussion in the schools.'' That is
substantially what happened, though not until he had left Rome: on
5~September 1607 Paul~V ended the congregation by permitting each
side to teach its own doctrine, forbidding either to censure the
other, and directing all to await a final decision of the
Apostolic See---a decision that was never issued; a 1611 decree
then restricted new publications on efficacious grace (CE 1908,
``Congregatio de Auxiliis''; K/E, reported). The non-resolution is
the historically important fact: the Holy See, at the height of
its confidence in deciding controversies, deliberately declined to
decide this one, on the line of Bellarmine's advice. Whether his
advice was principled restraint, party prudence protecting the
Jesuit position from censure, or both at once, the documents do
not say; the question is left open.
